% Part: first-order-logic
% Chapter: models-theories
% Section: set-theory

\documentclass[../../../include/open-logic-section]{subfiles}

\begin{document}

\olfileid{fol}{mat}{set}

\olsection{Een theorie over verzamelingen (verzamelingenleer)}

Bijna alle wiskunde kun je opbouwen binnen een theorie over
verzamelingen. Daarvoor zijn een paar dingen nodig. Om te beginnen
kiezen we axioma's voor de relatie~$\in$. Er bestaan verschillende
stelsels van zulke axioma's. Soms geven die stelsels zelfs
tegenstrijdige eigenschappen aan~$\in$. Eén stelsel wordt veruit het
meest gebruikt en bestudeerd: $\Log{ZFC}$, de
Zermelo--Fraenkel-verzamelingenleer met het keuzeaxioma. De axioma's van
dit stelsel zijn genoeg om bijna alles te bewijzen wat wiskundigen
verwachten te kunnen bewijzen. Maar eerst moeten we laten zien dat de
taal überhaupt kan \emph{zeggen} wat wiskundigen willen zeggen. De
taal heeft bijvoorbeeld geen !!{constant}s of !!{function}s. Daardoor
lijkt het eerst alsof we niet over één bepaalde verzameling kunnen
praten, zoals $\emptyset$ of $\Nat$. Ook bewerkingen op verzamelingen,
zoals $X \cup Y$ en $\Pow{X}$, lijken lastig uit te drukken. En dan
hebben we het nog niet eens over andere constructies, zoals relaties en
functies.

Naast ``is een element van'' gebruiken we ook de relatie ``is een
deelverzameling van''. Die tweede relatie is bijna even belangrijk. We
hoeven haar niet als nieuw basisteken aan de taal toe te voegen, want
we kunnen haar met ``is een element van'' \emph{definiëren}. Daarvoor
zoeken we in de taal van de verzamelingenleer !!a{formula}~$!A(x, y)$.
Een paar verzamelingen~$\tuple{X, Y}$ moet precies dan aan die formule
voldoen als $X \subseteq Y$. Wat betekent dat? $X$ is een
deelverzameling van $Y$ als ieder !!{element} van~$X$ ook
!!a{element} van~$Y$ is. Daarom definiëren we $\subseteq$ met de
formule
\[
\lforall[z][(z \in x \lif z \in y)]
\]
Vanaf nu kunnen we $\subseteq$ in iedere formule vervangen door deze
langere formule. We vervangen dan $x$ en $y$ op de juiste manier en
geven de gebonden variabele~$z$ zo nodig een andere naam. Een voorbeeld
is de extensionaliteit van verzamelingen. Die zegt hier: als de
verzamelingen~$x$ en $y$ elk een deelverzameling van de andere zijn,
dan zijn $x$ en $y$ dezelfde verzameling. We kunnen dit uitdrukken door
$\lforall[x][\lforall[y][((x \subseteq y \land y
    \subseteq x) \lif x = y)]]$. Vervangen we $\subseteq$ door de
definitie hierboven, dan krijgen we
\[
\lforall[x][\lforall[y][((\lforall[z][(z \in x \lif z \in y)] \land
    \lforall[z][(z \in y \lif z \in x)]) \lif x = y)]].
\]
Dit is inderdaad een van de axioma's van $\Log{ZFC}$. Het heet het
``extensionaliteitsaxioma''.

De taal heeft geen !!{constant} die $\emptyset$ aanwijst. Toch kunnen
we zeggen dat ``$x$ leeg is'': we schrijven $\lnot \lexists[y][y \in
x]$, dus dat er geen enkel element in die verzameling zit. De bewering ``$\emptyset$
bestaat'' wordt daarmee de !!{sentence}~$\lexists[x][\lnot
\lexists[y][y \in x]]$. Ook dit is een axioma van~$\Log{ZFC}$.
(Uit het extensionaliteitsaxioma volgt bovendien dat er maar één lege
verzameling is.) Als we in deze taal over $\emptyset$ willen praten,
schrijven we dus steeds iets als ``er is een verzameling die leeg is en
\dots'' Als voorbeeld kunnen we zeggen dat $\emptyset$ een
deelverzameling van iedere verzameling is. Dan schrijven we
\[
\lexists[x][(\lnot\lexists[y][y \in x] \land \lforall[z][x \subseteq
    z])]
\]
Daarbij moeten we $x \subseteq z$ natuurlijk weer vervangen door de
definitie daarvan.

Voor bewerkingen op verzamelingen, zoals $X \cup Y$ en $\Pow{X}$,
gebruiken we dezelfde aanpak. De taal van de verzamelingenleer heeft
geen functiesymbolen. We beschrijven daarom niet een bewerkingsteken,
maar de functionele relaties $X \cup Y = Z$ en $\Pow{X} = Y$. Dat kan
met
\begin{align*}
& \lforall[u][((u \in x \lor u \in y) \liff u \in z)]\\
& \lforall[u][(u \subseteq x \liff u \in y )]
\end{align*}
De eerste formule zegt dat de !!{element}s van $X \cup Y$ precies de
verzamelingen zijn die in~$X$ of in~$Y$ zitten. De tweede zegt dat de
!!{element}s van $\Pow{X}$ precies de deelverzamelingen van~$X$ zijn.
Maar hierdoor mogen we $x \cup y$ of $\Pow{x}$ nog niet gebruiken alsof
het termen zijn. We mogen alleen de volledige !!{formula}s gebruiken
die de relaties $X \cup Y = Z$ en $\Pow{X} = Y$ definiëren. Ook weten
we nog niet of er altijd een verzameling bestaat die aan zo'n formule
voldoet. We weten dus nog niet of iedere vereniging en iedere
machtsverzameling bestaat. De !!{sentence}
$\lforall[x][\lexists[y][\Pow{x} = y]]$ legt dat bestaan voor
machtsverzamelingen vast. Dit is een ander axioma van~$\Log{ZFC}$ (het
machtsverzamelingsaxioma).

Hoe praten we dan over geordende paren en functies? We moeten laten
zien hoe we die als bijzondere soorten verzamelingen kunnen opbouwen.
Een mogelijkheid is om het geordende paar $\tuple{x, y}$ te definiëren
als de verzameling $\{\{x\}, \{x, y\}\}$. Ook nu kunnen we geen
!!a{function} invoeren die deze verzameling aanwijst. We kunnen alleen
de relatie $\tuple{x, y} = z$ definiëren. Die relatie zegt hier precies
dat $\{\{x\}, \{x, y\}\} = z$:
\[
\lforall[u][(u \in z \liff (\lforall[v][(v \in u \liff v = x)] \lor
  \lforall[v][(v \in u \liff (v = x \lor v = y))]))]
\]
Dit zegt het volgende. De !!{element}s~$u$ van~$z$ zijn precies de
verzamelingen die alleen $x$ als !!{element} hebben, of alleen $x$
en~$y$ als !!{element}s. Het zijn dus precies de verzamelingen die
gelijk zijn aan $\{x\}$ of aan $\{x, y\}$. Met deze beschrijving kunnen
we verdergaan. We kunnen bijvoorbeeld zeggen dat $X \times Y = Z$:
\[
\lforall[z][(z \in Z \liff \lexists[x][\lexists[y][(x \in
        X \land y \in Y \land \tuple{x, y} = z)]])]
\]

We kunnen een functie $f \colon X \to Y$ nu coderen als de relatie
$f(x) = y$. De functie wordt dan de verzameling
paren~$\Setabs{\tuple{x,y}}{f(x) = y}$. Een verzameling~$f$ is een
functie van $X$ naar $Y$ als aan drie voorwaarden is voldaan. Ten eerste
is zij een relatie $\subseteq X \times Y$. Ten tweede krijgt iedere
invoer $x \in X$ minstens één uitkomst $y \in Y$ waarvoor
$\tuple{x, y} \in f$; de relatie is dan totaal. Ten derde krijgt
een invoer nooit twee verschillende uitkomsten: als $\tuple{x, y},
\tuple{x, y'} \in f$, dan geldt $y = y'$. De relatie is dan
functioneel, want de waarde van een functie moet uniek zijn. We
kunnen ``$f$ is een functie van $X$ naar $Y$'' dus schrijven als:
\begin{align*}
 \lforall[u][(u \in f \lif {}] & \lexists[x][\lexists[y][(x \in X \land y \in
      Y \land \tuple{x, y} = u)]]) \land {}\\
 \lforall[x][(x \in X \lif {}] &
      (\lexists[y][(y \in Y \land \mathrm{maps}(f, x, y))] \land {}\\
& (\lforall[y][\lforall[y'][((\mathrm{maps}(f, x, y) \land
    \mathrm{maps}(f, x, y')) \lif y = y')]]))
\end{align*}
Hier is $\mathrm{maps}(f, x, y)$ een afkorting voor
$\lexists[v][(v \in f \land \tuple{x, y} = v)]$. Deze !!{formula}
zegt dat ``$f(x) = y$''.

We kunnen nu ook uitdrukken dat $f\colon X \to Y$ !!{injective} is.
Een mogelijke formule is:
\begin{multline*}
 f \colon X \to Y \land \lforall[x][\lforall[x'][((x \in X \land x' \in
     X \land {}]] \\
 \lexists[y][(\mathrm{maps}(f, x, y) \land \mathrm{maps}(f,
        x', y))]) \lif x = x')
\end{multline*}
Een functie~$f\colon X \to Y$ is !!{injective} als twee invoeren alleen
dezelfde uitkomst kunnen hebben wanneer ze gelijk zijn. Dus: als $f$
de elementen $x, x' \in X$ allebei naar één~$y$ stuurt, dan geldt $x =
x'$. Kort de formule hierboven af als $\mathrm{inj}(f, X, Y)$. Dan
kunnen we in de taal van de verzamelingenleer al de stelling van Cantor
opschrijven: er bestaat geen !!{injective} functie van $\Pow{X}$ naar
$X$:
\[
\lforall[X][\lforall[Y][(\Pow{X} = Y \lif
    \lnot\lexists[f][\mathrm{inj}(f, Y, X)])]]
\]

Misschien lijkt er nog een axioma nodig dat bij iedere eigenschap een
verzameling laat bestaan. Neem een formule $!A(x)$ van de
verzamelingenleer waarin de variabele~$x$ vrij voorkomt. Dan kunnen we
de volgende !!{sentence} bekijken:
\[
\lexists[y][\lforall[x][(x \in y \liff !A(x))]].
\]
Deze !!{sentence} zegt dat er een verzameling~$y$ bestaat. De
!!{element}s daarvan zijn precies de $x$ die aan~$!A(x)$ voldoen. Dit
vaste patroon heet het ``comprehensieprincipe''. Het ziet er nuttig uit,
maar leidt helaas tot een tegenspraak. Neem $!A(x) \ident \lnot x \in
x$. Het comprehensieprincipe zegt dan
\[
\lexists[y][\lforall[x][(x \in y \liff x \notin x)]],
\]
Het eist dus een verzameling van alle verzamelingen die geen
!!{element}s van zichzelf zijn. Zo'n verzameling kan niet bestaan. Dit
is de paradox van Russell. $\Log{ZFC}$ gebruikt daarom een
beperkte---en consistente---versie van het principe: het
separatieprincipe. Dat principe begint met een verzameling die al
bestaat en neemt daaruit alleen de elementen die aan de gekozen
eigenschap voldoen:
\[
\lforall[z][\lexists[y][\lforall[x][(x \in y \liff (x \in z \land
      !A(x))]]].
\]

\begin{prob}
Laat met !!a{derivation} zien dat het comprehensieprincipe inconsistent
is. Leid daarvoor af dat
\[
\lexists[y][\lforall[x][(x \in y \liff x \notin x)]] \Proves \lfalse.
\]
Het kan helpen om eerst $(A \lif \lnot A) \land (\lnot A \lif A)
\Proves \lfalse$ af te leiden.
\end{prob}
\end{document}
